\section{System Model}
\label{sec:model}

\begin{table}[!t]
\caption{Simulation parameters. Every value is read from the single frozen configuration object used by all ten protocols; nothing is protocol-specific.}
\label{tab:parameters}
\centering
\begin{tabular}{ll}
\toprule
Parameter & Value \\
\midrule
\multicolumn{2}{l}{\textit{Topology and energy}} \\
\quad Nodes $N$ & 100 \\
\quad Field & 100 $\times$ 100 m \\
\quad Base station & $(50, 150)$ m \\
\quad Initial energy $E_0$ & 1.0 J per node \\
\quad Rounds (cap) & 7,000 \\
\quad Runs per configuration & 30 \\
\addlinespace
\multicolumn{2}{l}{\textit{Packets}} \\
\quad Data packet & 4,000 bits \\
\quad Control packet & 200 bits \\
\addlinespace
\multicolumn{2}{l}{\textit{Radio (first-order model)}} \\
\quad $E_{\mathrm{elec}}$ & 50 nJ/bit \\
\quad $\varepsilon_{\mathrm{fs}}$ & 10 pJ/bit/m$^2$ \\
\quad $\varepsilon_{\mathrm{mp}}$ & 0.0013 pJ/bit/m$^4$ \\
\quad $E_{\mathrm{DA}}$ (fusion) & 5 nJ/bit \\
\quad Crossover $d_0$ & 87.7 m \\
\addlinespace
\multicolumn{2}{l}{\textit{Channel}} \\
\quad Transmit power $P_t$ & 12 dBm \\
\quad Path loss at 1 m & 40 dB \\
\quad Path-loss exponent $n$ & 3.0 \\
\quad Shadowing $\sigma$ & 4 dB \\
\quad Noise floor & -105 dBm \\
\quad ARQ retransmissions & 2 \\
\addlinespace
\multicolumn{2}{l}{\textit{Sensed data (Ornstein--Uhlenbeck)}} \\
\quad Mean $\mu$, reversion $\theta$ & 50, 0.05 \\
\quad Step $\sigma$ & 2.0 \\
\quad TEEN/APTEEN $H_T$, $S_T$, $T_C$ & 55, 1.0, 50 rounds \\
\addlinespace
\multicolumn{2}{l}{\textit{Protocol-agnostic}} \\
\quad Head fraction $p$ & 0.05 \\
\quad Density radius & 25 m \\
\quad NSGA-II population, generations & 40, 30 \\
\quad Re-clustering interval & 5 rounds \\
\addlinespace
\multicolumn{2}{l}{\textit{Timing}} \\
\quad Bit rate & 250 kbps \\
\quad Slot / guard / processing & 20 / 4 / 1 ms \\
\bottomrule
\end{tabular}
\end{table}


\subsection{Network and Energy}

One hundred nodes are placed uniformly at random in a $100\times100$~m field,
with the base station outside the field at $(50, 150)$, so node-to-sink
distances span 50 to 158~m. Every node begins with $E_0 = 1$~J, is stationary,
and fails only by depletion. The configuration carries fields for two-level and
three-level heterogeneity~\cite{smaragdakis2004sep,qing2006deec}, disabled
throughout, so every result here is for a homogeneous population.
Table~\ref{tab:parameters} lists every parameter; nothing in it is
protocol-specific, and no constant is defined outside the single frozen
configuration object.

A run proceeds in rounds until the last node dies or a cap of 7{,}000 rounds is
reached. Energy is a linear reservoir with a single charging function: a node
asked to spend more than it holds spends what remains, dies, and loses the
packet in flight. Refusing the transmission instead would strand the remainder
and break the conservation identity
\begin{equation}
\sum_i \left( E_0(i) - E_{\mathrm{res}}(i) \right)
= E_{\mathrm{tx}} + E_{\mathrm{rx}} + E_{\mathrm{ctrl}}
+ E_{\mathrm{agg}} + E_{\mathrm{retry}},
\label{eq:conservation}
\end{equation}
asserted at the end of every run to $10^{-9}$~J and the primary correctness
check on the simulator.

\subsection{Radio Energy}

We use the first-order model of Heinzelman et
al.~\cite{heinzelman2000leach,heinzelman2002appspecific}. Transmitting $k$ bits
over distance $d$ costs
\begin{equation}
E_{\mathrm{tx}}(k, d) =
\begin{cases}
k E_{\mathrm{elec}} + k \varepsilon_{\mathrm{fs}} d^{2}, & d < d_0 \\[2pt]
k E_{\mathrm{elec}} + k \varepsilon_{\mathrm{mp}} d^{4}, & d \geq d_0
\end{cases}
\label{eq:radio}
\end{equation}
with $E_{\mathrm{rx}}(k) = k E_{\mathrm{elec}}$ and a fusion charge of
$E_{\mathrm{DA}}$ per bit received. The crossover
$d_0 = \sqrt{\varepsilon_{\mathrm{fs}}/\varepsilon_{\mathrm{mp}}} = 87.7$~m sits
inside the range of node-to-sink distances, which makes head placement relative
to $d_0$ a live variable: a head at 100~m pays the fourth-power branch and one
at 80~m does not. In the seed-0 topology 36 of the 100 nodes fall inside $d_0$,
against an expected 32.7\% over the placement distribution.

The model is minimal and its omissions matter. There is no idle listening, no
sleep-wake transition energy, no radio startup and no processor energy. Idle
listening often dominates real budgets~\cite{ye2002smac,polastre2004bmac}, and
including it would compress the differences between all nine protocols, so this
bounds every energy figure that follows.

\subsection{Channel}
\label{sec:channel}

Packet loss is derived from a link budget rather than assigned as a constant.
Path loss follows a log-distance law with per-link
shadowing~\cite{rappaport2002wireless},
\begin{equation}
PL(d) = PL_{1\mathrm{m}} + 10 n \log_{10} d + X_\sigma,
\qquad X_\sigma \sim \mathcal{N}(0, \sigma^2),
\label{eq:pathloss}
\end{equation}
giving a receive power of $P_t - PL(d)$ and an SNR against a $-105$~dBm noise
floor. Bit errors follow the closed-form non-coherent BFSK
expression~\cite{proakis2008digital}
$\mathrm{BER} = \tfrac{1}{2}e^{-\gamma/2}$, and
$\mathrm{PER} = 1 - (1-\mathrm{BER})^{k}$, evaluated as
$-\mathrm{expm1}\!\left(k \ln(1-\mathrm{BER})\right)$ for numerical stability.

\begin{figure}[!t]
\centering
\includegraphics[width=\columnwidth]{figA_per_curve.png}
\caption{Packet error rate against distance for a 4{,}000-bit packet, with the
crossover $d_0 = 87.7$~m and the measured band of node-to-sink distances marked.
The curve is a waterfall, not a gradient: effectively zero out to about 120~m,
saturated by 158~m. Where a protocol places the far tail of its head-to-sink hop
relative to that step decides what it pays in retransmissions, which is why
transmit power is an experimental factor in Section~\ref{sec:design} rather than
a constant.}
\label{fig:per}
\end{figure}

Fig.~\ref{fig:per} plots the result, and it is a waterfall rather than a
gradient. Without shadowing, PER is $2.6\times10^{-8}$ at 100~m,
$1.0\times10^{-3}$ at 120~m, 0.19 at 140~m and 0.97 at 158~m, so the transition
from usable to useless spans about 3.6~dB of SNR and 40~m. Shadowing at
$\sigma = 4$~dB smears that step into a distribution over links and is the
source of the genuine per-topology variance in delivery ratio; it is drawn once
per topology and held fixed, because the nodes are static. Failed packets are
retried up to twice, each attempt paying full transmit and receive energy, which
is what makes a lossy channel cost joules rather than only packets. Retry energy
is charged to a separate category so it can be measured directly rather than
inferred.

\subsection{Control-Traffic Accounting}
\label{sec:control}

Control traffic is charged by the engine according to the protocol's declared
class, not by the protocol itself, and both endpoints pay. Distributed protocols
pay a full handshake every round: an advertisement priced at the distance to the
farthest joiner with a receive charge on every living non-head, a join request
per member, and a TDMA schedule broadcast every member receives. The chain
protocol passes a leader token hop by hop instead, plus a sink broadcast when
the chain is rebuilt. Centralized protocols skip the advertisement and
join-request steps, since the sink has already told each node its assignment,
but do still pay the per-round schedule broadcast under identical pricing. Their
uplink node state is treated as piggybacked on data traffic and charged nothing.
The subsidy is therefore narrower than ``centralized protocols pay almost no
control traffic'': it is specifically the advertisement, the join request and
the uplink report, which is why the measured control-energy share of the
centralized five is 2.33\% to 3.95\% rather than zero.

That choice is the largest modeling decision in the paper.
Section~\ref{sec:subsidy} gives its magnitude and what it costs every
cross-class comparison.

\subsection{Sensed Data, Latency and Invariants}
\label{sec:invariants}

Each node's reading follows a discrete Ornstein-Uhlenbeck
process~\cite{uhlenbeck1930brownian} stationary at approximately
$\mathcal{N}(50, 6.4^2)$; a pure random walk was rejected because over 7{,}000
rounds it pegs against any clipping bound, which would make the reactive
protocols' behavior an artifact of the clip. TEEN and APTEEN's thresholds are
set against that stationary distribution, hard at $+0.78\sigma$ and soft at
$0.16\sigma$, and were not adjusted after observing any result. The full sensed
stream is generated once per topology, so every protocol reads byte-identical
values.

Latency is analytic,
$L = s\,T_{\mathrm{slot}} + h\,(T_{\mathrm{tx}} + T_{\mathrm{proc}})$, with $s$
the sender's TDMA slot index and $h$ its hop depth. A head's own packet carries
a slot index equal to its member count, which is what makes cluster latency
scale with cluster size. There is no queueing or contention delay, so absolute
figures are lower bounds and only the ordering should be read.

Three rules cannot be bypassed by a protocol implementation. Transmissions
execute sorted by hop depth and then node index, because plain index ordering
breaks causality for roughly half of all head-member pairings on a random
topology and silently corrupts fusion totals. A node that receives data emits at
most one outgoing packet and pays the fusion charge on every bit received. And
the count of original readings in a delivered packet is computed by the engine,
not declared by the protocol, since an earlier version that let protocols
declare it would have let a protocol inflate its own delivery figures.
